\section{Expériences professionnelles}
\resumeSubHeadingListStart

    \resumeSubheading
    {Stagiaire de recherche}{2025 - 2026}
    {Recherche théorique et pratique en Quantum Error Mitigation (QEM)}{Qiskit, SciPy, NumPy}
    {IBM Quantum}
        \resumeItemListStart
            \resumeItem{Recherche sur une nouvelle technique de QEM mêlant théorie quantique et data science}
            \resumeItem{Collaboration avec des chercheurs d'IBM pour implémenter et tester la technique sur des processeurs quantiques réels}
            \resumeItem{Exploration de méthodes de synthèse de circuits basées sur des diagrammes de décision pour les circuits dynamiques}
        \resumeItemListEnd

        \resumeSubheading
    {Ingénieur logiciel}{2025}
    {Césure très sélective : trois projets innovants de 7 semaines pour des entreprises}{Rust, C++, Flutter, Docker\ghicon{fhe-bpce}}
    {Paris Digital Lab}
        \resumeItemListStart
            \resumeItem{Programmation d'un framework Cloud-native performant de Fully Homomorphic Encryption pour le Groupe BPCE}
            \resumeItem{Développement d'une application sécurisée de téléphonie et messagerie avec vérification d'identité eIDAS / EUDI}
        \resumeItemListEnd

    \resumeSubheading
    {Stagiaire de recherche en Informatique Quantique}{2023 - 2025}
        {Stagiaire à temps partiel au sein de l'équipe QuaCS de l'INRIA}{QASM, C++, GoogleTest, CMake\ghicon{coto}}
        {INRIA}
        \resumeItemListStart
            \resumeItem{Théorisation d'une structure de données innovante pour représenter des états et circuits quantiques}
            \resumeItem{Commencement de l'écriture d'un article sur le sujet, détaillant les propriétés mathématiques de la structure de données}
            \resumeItem{Création et implémentation d'algorithmes pour la simulation sur machines classiques de circuits quantiques}
        \resumeItemListEnd

    \resumeSubheading
    {Secrétaire Général \& Directeur des formations}{2023 - 2025}
        {Administrateur réseau \& systèmes du plus grand FAI étudiant de France}{Kubernetes, Ansible, OpenStack, Ceph}
        {ViaRézo (association IT)}
        \resumeItemListStart
        \resumeItem{Operation de 400+ bornes Wi-Fi Unifi et 50+ switches Juniper pour fournir internet à 2,5k+ étudiants}
        \resumeItem{Administration d'un cluster OpenStack de 200+ machines virtuelles et d'un cluster Kubernetes de 150+ pods}
            \resumeItem{Responsabilité légale des bases de données (13M+ lignes) et de leur conformité (RGPD, NIS2)}
            \resumeItem{Conception et direction de formations et ateliers techniques (Bash, Docker, \LaTeX, Git)}
        \resumeItemListEnd

    \resumeSubheading
    {Stagiaire en infrastructure réseau}{2024}
        {\resumeItemBackwards{Dépannage et sécurisation de l'infrastructure réseau d'une bibliothèque universitaire}}{Cisco IOS, Wireshark, VLANs}
        {Université Paris-Saclay}

\resumeSubHeadingListEnd
